\def\baselinestretch{1}
\chapter{Research Agenda and Acknowledgement}
\ifpdf
    \graphicspath{{Conclusions/ConclusionsFigs/PNG/}{Conclusions/ConclusionsFigs/PDF/}{Conclusions/ConclusionsFigs/}}
\else
    \graphicspath{{Conclusions/ConclusionsFigs/EPS/}{Conclusions/ConclusionsFigs/}}
\fi

\def\baselinestretch{1.66}

\section{Research Agenda}
We summarize our future research plan and agenda in \cref{tab:agenda}.
\begin{table}[H]
    \centering
    \begin{tabular}{|l|p{8cm}|}
        \hline
        Period & Research Plan \\
        \hline
        2024.7-2024.12 & SincKAN: idea formulation, model construction, numerical experiment, paper writing, submission to conference, peer review and rebuttal. \\
        \hline
        2024.12-2025.5 & Continue submission and rebuttal process of SincKAN; Continue project on machine learning for phase-field model \\
        \hline
        2025.5-2025.12 & Conduct extension work on SincKAN; Paper submission for phase-field model project \\
        \hline
        2026.1-2026.12 & Extension work on phase-field model project; Other AI for multiphysics modeling projects \\
        \hline
    \end{tabular}
    \caption{Thesis research agenda}
    \label{tab:agenda}
\end{table}
\section{Acknowledgement}
Beneath the lush lychee trees and amidst the hazy morning mist, spring in the southern land quietly ascends the green branches once more. Unknowingly, I have already spent three years at the Southern University of Science and Technology. Among these years, what I cherish the most are the earnest teachings of my advisor, Professor Yang Jiang, who has significantly guided my rapid growth in the interdisciplinary field of computational mathematics and machine learning. In Professor Yang’s group, we enjoy ample room for exploration and growth, allowing us to lay a solid foundation in the traditional domain of computational mathematics while venturing into the rapidly evolving field of machine learning, exploring new algorithms and applications. This has been profoundly beneficial for both advancing science and personal career development. Professor Yang provided invaluable guidance during the early stages of my doctoral research, which included projects such as the gradient ascent algorithm for adaptive sampling in PINNs, the application of SincKAN to fractional partial differential equations, and the use of machine learning for approximating phase-field free energies. Without his support, many of the results presented in this report would not have been possible.  

I would also like to extend my sincere gratitude to my collaborator Mr. Yu Tianchi, who is currently pursuing his PhD at the Skolkovo Institute of Science and Technology in Moscow. He made critical contributions to the experimental work on residual gradient methods for Physics-Informed Neural Networks (PINNs) and helped me gain a deeper understanding of Sinc numerical methods. His profound expertise in scientific computing and machine learning, as well as his exceptional technical skills, have deeply impressed me and inspired me to further enhance my own skills to address technical shortcomings.



% ------------------------------------------------------------------------

%%% Local Variables: 
%%% mode: latex
%%% TeX-master: "../thesis"
%%% End: 




